\chapter{Arterial Line Placement and Arterial Blood Gas Sampling}

\section*{Overview}
Arterial line placement is the insertion of a small catheter into a peripheral or central artery for continuous blood pressure monitoring and repeated blood sampling. Arterial blood gas (ABG) sampling is the withdrawal of arterial blood to assess oxygenation, ventilation, and acid-base status.

\section*{Alternative Names}
A-line, arterial cannulation, radial artery catheterization, ABG sampling

\section*{Principle}
An arterial catheter transduces the phasic pressure waveform of arterial blood to a monitor, providing beat-to-beat blood pressure and enabling repeated sampling without repeated needle sticks. ABG analysis measures pH, partial pressures of oxygen and carbon dioxide, and derived parameters such as bicarbonate and base excess, reflecting gas exchange and metabolic status.

\section*{History}
Arterial cannulation was described in the 20th century and became widespread in intensive care during the mid-1900s. The Allen test was historically used to assess collateral ulnar circulation before radial cannulation, though its predictive value is debated.

\section*{Why This Test or Procedure Is Ordered}
Ordered for continuous blood pressure monitoring in hemodynamically unstable patients, frequent ABG or lab sampling, and titration of vasoactive drugs or mechanical ventilation. ABG sampling is obtained to diagnose respiratory failure, evaluate acid-base disorders, and guide oxygenation and ventilation targets.

\section*{Is This a Primary or Secondary Test or Procedure}
A primary monitoring and diagnostic procedure in critically ill patients requiring real-time hemodynamic or blood gas assessment.

\section*{How This Test or Procedure Is Performed}
The artery, most often the radial, is palpated or visualized with ultrasound, and the overlying skin is prepared and anesthetized. The catheter is advanced over a guidewire or directly into the artery until pulsatile blood return is seen, then secured and connected to a pressurized, heparinized transducer system. For isolated ABG sampling, a small volume of blood is drawn anaerobically with a heparinized syringe and processed immediately.

\section*{What Preparation Is Needed}
Informed consent, assessment of distal perfusion and collateral flow, and sterile technique. The Allen test or ultrasound-guided confirmation of vessel patency is often performed before radial cannulation.

\section*{Contraindications and Precautions}
Relative contraindications include inadequate collateral circulation, Raynaud phenomenon, peripheral vascular disease, coagulopathy, and local infection. Precautions include avoiding overly aggressive blood sampling, maintaining sterile technique, and prompt removal when the catheter is no longer needed.

\section*{Risks}
Hematoma, bleeding, arterial thrombosis with distal ischemia, infection, pseudoaneurysm, and accidental intra-arterial injection of medications.

\section*{What Happens After the Test or Procedure}
The line is kept patent with a continuous flush solution and secured to prevent dislodgement; the site is inspected regularly for ischemia or infection. The catheter is removed with firm pressure when monitoring is no longer required, and the site is observed for bleeding.

\section*{Understanding Your Results}
ABG results show oxygenation (PaO2), ventilation (PaCO2), and pH; low PaO2 with normal PaCO2 suggests hypoxemic respiratory failure, while elevated PaCO2 with acidosis indicates ventilatory failure. Continuous arterial pressure readings guide vasopressor and fluid therapy in shock states.

\section*{Normal Results}
Normal ABG values include a pH of 7.35--7.45, PaO2 of 80--100 mm Hg on room air, and PaCO2 of 35--45 mm Hg, with mean arterial pressure near 70--100 mm Hg.

\section*{Follow-up Tests and Procedures}
Repeat ABG and other labs are drawn through the line as indicated by the clinical course. When the line is removed, assessment of distal perfusion and, if clinically indicated, pulse oximetry or vascular ultrasound follow.

\section*{References}
\begin{enumerate}
  \item MedlinePlus. Blood gases. U.S. National Library of Medicine; 2025.
  \item Current Medical Diagnosis and Treatment. McGraw-Hill; 2025.
\end{enumerate}

\clearpage
